The experimental results are evaluated using 
the following metrics:

\subsubsection{Classification Performance}

The model's classification performance is 
measured using top-1 accuracy on the CIFAR-10 
test set:

\begin{equation}
\text{Accuracy} = \frac{1}{n} \sum_{i=1}^{n} 
\mathbf{1}[\hat{y}_i = y_i]
\label{eq:accuracy}
\end{equation}

where $\mathbf{1}[\cdot]$ is the indicator 
function and $n$ is the number of test examples.

\subsubsection{Adversarial Robustness}

The model's robustness to adversarial attacks 
is measured by the drop in accuracy under 
FGSM and PGD perturbations:

\begin{equation}
\Delta_{adv} = \text{Accuracy}(f_\theta, 
\mathcal{D}_{test}) - \text{Accuracy}
(f_\theta, \mathcal{D}_{adv})
\label{eq:advdrop}
\end{equation}

where $\mathcal{D}_{adv}$ denotes the 
adversarially perturbed test set. A larger 
$\Delta_{adv}$ indicates lower robustness 
and higher vulnerability to adversarial attacks.

\subsubsection{Memorization Rate}

The memorization rate is defined as the 
proportion of training examples classified 
as memorized:

\begin{equation}
\text{MR} = \frac{|\mathcal{M}|}{|\mathcal{D}|}
\label{eq:memrate}
\end{equation}

where $|\mathcal{M}|$ is the number of memorized 
examples and $|\mathcal{D}|$ is the total number 
of training examples.

\subsubsection{Attribution Similarity}

To quantify the difference between attribution 
maps of memorized and non-memorized examples, 
the structural similarity index (SSIM) is used:

\begin{equation}
\text{SSIM}(\mathcal{A}_1, \mathcal{A}_2) = 
\frac{(2\mu_1\mu_2 + c_1)(2\sigma_{12} + c_2)}
{(\mu_1^2 + \mu_2^2 + c_1)
(\sigma_1^2 + \sigma_2^2 + c_2)}
\label{eq:ssim}
\end{equation}

where $\mu_1, \mu_2$ are the mean values of 
the attribution maps, $\sigma_1^2, \sigma_2^2$ 
are their variances, $\sigma_{12}$ is their 
covariance, and $c_1, c_2$ are stability 
constants. A low SSIM score between memorized 
and non-memorized attribution maps indicates 
that the model relies on fundamentally different 
features for these two groups of examples.